\newpage

\section{নাহি}

\noindent যে ফেল দিয়ে কাজ থেকে নিষেধ বোঝায়, তাকে
{\arabicfont \textarabic{فعل نهي}} বলে।
যেমন {\arabicfont \textarabic{لَا تَفْعَلْ}} করো না,
{\arabicfont \textarabic{لَا تَكْذِبْ}} মিথ্যা বলো না।

\vspace{0.3cm}

\noindent নাহি মুজারে মারুফ থেকে হয়।

\begin{enumerate}
  \item শুরুতে {\arabicfont \textarabic{\sfx{لَا}}} বসাও।
  \item পাঁচ সিগায় শেষে সাকিন:
        গায়েব পুং একবচন, গায়েব স্ত্রী একবচন,
        হাজির পুং একবচন, মুতাকাল্লিমের দুই সিগা।
        লাম ইল্লাত হলে লাম ফেলো।
        যেমন {\arabicfont \textarabic{لَا تَرْمِ، لَا تَخْشَ}}।
  \item সাত সিগা থেকে নুনে ই'রাব ফেলো:
        দ্বিবচন চার, পুং বহুবচন দুই, হাজির স্ত্রী একবচন।
  \item স্ত্রী বহুবচনের দুই সিগা আগের মতো থাকে।
\end{enumerate}

\vspace{0.15cm}

\begin{center}
  {\arabicfont \color{titlegreen} \textbf{\textarabic{نهي حاضر معروف}}}
\end{center}

\vspace{0.15cm}

\begingroup
\setlength{\arrayrulewidth}{0.5pt}
\setlength{\tabcolsep}{3pt}
\renewcommand{\arraystretch}{1.9}

\noindent\begin{tabularx}{\textwidth}{|C|C|}
  \hline
  \rowcolor{headergreen}
  {\color{white}\textbf{অর্থ}} &
  {\color{white}\textbf{রূপ}} \\
  \hline
  \rowcolor{lightgreen}
  তুমি (একজন পুরুষ) করো না & {\arabicfont \nahi} \\
  \hline
  তোমরা (দুজন পুরুষ) করো না & {\arabicfont \nahiaa} \\
  \hline
  \rowcolor{lightgreen}
  তোমরা (সকল পুরুষ) করো না & {\arabicfont \nahiuu} \\
  \hline
  তুমি (একজন স্ত্রী) করো না & {\arabicfont \nahii} \\
  \hline
  \rowcolor{lightgreen}
  তোমরা (দুজন স্ত্রী) করো না & {\arabicfont \nahiaa} \\
  \hline
  তোমরা (সকল স্ত্রী) করো না & {\arabicfont \nahina} \\
  \hline
\end{tabularx}
\endgroup

\vspace{0.25cm}

\noindent গায়েব ও মুতাকাল্লিমেও শুরুতে {\arabicfont \textarabic{\sfx{لَا}}}
এবং শেষে সেই জযম।
যেমন {\arabicfont \ovl{\mk{لَا يَ\zwj}\rt{\zwj فْعَل\nh{ْ}}}}।

\vspace{0.15cm}

\begin{center}
  {\small
    লাল = {\arabicfont \textarabic{لَا}}, হরফে মুজারা ও সিগার শেষ
    \enspace$\bullet$\enspace
    কালো = মূল
  }
\end{center}
